\subsection{Lee--Yang 四重覆盖的自纤维积几何、Prym 因子与有限域二阶矩}\label{subsec:conclusion-leyang-self-fiber-product-prym-second-moment}
本小节在既定四次谱曲线
\[
\Pi(\lambda,y)=\lambda^{4}-\lambda^{3}-(2y+1)\lambda^{2}+\lambda+y(y+1)=0
\]
及其椭圆化覆盖 \(\pi_y:E\to\PP^1_y\) 的框架中，给出非对角自纤维积分量的显式方程、交换商曲线的超椭圆模型、有限域纤维二阶矩恒等式、Jacobian 同源分解以及属 \(2\) 商模型的分歧判别式。

\begin{theorem}[非对角自纤维积分量的显式模型与属数]\label{thm:conclusion-leyang-self-fiber-product-genus6}
设 \(D\) 为
\[
E\times_{\PP^1_y}E
\]
去除对角分量后再归一化得到的平滑射影曲线。则：
\begin{enumerate}
  \item 在 \(\PP^1_{\lambda_1}\times\PP^1_{\lambda_2}\) 上，\(D\) 的非对角双有理模型由
  \[
  R(\lambda_1,\lambda_2)=0
  \]
  给出，其中
  \[
  \begin{aligned}
  R(\lambda_1,\lambda_2)=\;&\lambda_1^6+2\lambda_1^5\lambda_2-2\lambda_1^5-\lambda_1^4\lambda_2^2-4\lambda_1^4\lambda_2+\lambda_1^4\\
  &-4\lambda_1^3\lambda_2^3-2\lambda_1^3\lambda_2^2+2\lambda_1^3\lambda_2+2\lambda_1^3\\
  &-\lambda_1^2\lambda_2^4-2\lambda_1^2\lambda_2^3+3\lambda_1^2\lambda_2^2+6\lambda_1^2\lambda_2-3\lambda_1^2\\
  &+2\lambda_1\lambda_2^5-4\lambda_1\lambda_2^4+2\lambda_1\lambda_2^3+6\lambda_1\lambda_2^2-4\lambda_1\lambda_2\\
  &+\lambda_2^6-2\lambda_2^5+\lambda_2^4+2\lambda_2^3-3\lambda_2^2+1.
  \end{aligned}
  \]
  \item 映射 \(D\to\PP^1_y\) 为次数 \(12\) 的连通分歧覆盖，且
  \[
  g(D)=6.
  \]
  \item 两个自然投影 \(\operatorname{pr}_i:D\to E\)（\(i=1,2\)）均为次数 \(3\) 的有限态射。
\end{enumerate}
\end{theorem}
\begin{proof}
将 \(\Pi(\lambda,y)\) 视为关于 \(y\) 的二次式：
\[
\Pi(\lambda,y)=y^2+(1-2\lambda^2)y+(\lambda^4-\lambda^3-\lambda^2+\lambda).
\]
因此 \((\lambda_1,\lambda_2)\) 对应同一 \(y\) 当且仅当
\[
\operatorname{Res}_{y}\bigl(\Pi(\lambda_1,y),\Pi(\lambda_2,y)\bigr)=0.
\]
直接计算可得
\[
\operatorname{Res}_{y}\bigl(\Pi(\lambda_1,y),\Pi(\lambda_2,y)\bigr)
=(\lambda_1-\lambda_2)^2\,R(\lambda_1,\lambda_2).
\]
其中 \((\lambda_1-\lambda_2)^2\) 对应对角分量，余因子 \(R\) 即非对角分量方程；其归一化即 \(D\)。

关于次数与属数，沿用定理 \ref{thm:xi-terminal-zm-leyang-elliptic-structure} 中 \(\mathrm{Mon}(\pi_y)\cong S_4\) 与分支护照 \((2,2,2,2,2,4)\)。
非对角有序对集
\[
\Omega:=\{(i,j)\in\{1,2,3,4\}^2:\ i\neq j\},
\qquad |\Omega|=12,
\]
上作用传递，故 \(D\to\PP^1_y\) 次数为 \(12\) 且连通。
有限简单分歧处局部单值为换位，其在 \(\Omega\) 上循环型为 \(2^5\,1^2\)，贡献 \(5\)；
无穷远处局部单值为四循环，在 \(\Omega\) 上循环型为 \(4^3\)，贡献 \(3\cdot(4-1)=9\)。
总分歧度
\[
5\cdot 5+9=34.
\]
由 Riemann--Hurwitz：
\[
2g(D)-2=-2\cdot 12+34=10,
\]
故 \(g(D)=6\)。

对一般 \(P\in E\)（其 \(y(P)\) 非分歧），纤维 \(\pi_y^{-1}(y(P))\) 有 \(4\) 个互异点；
固定第一坐标为 \(P\) 时，第二坐标可取其余 \(3\) 点，故 \(\deg(\operatorname{pr}_1)=3\)，同理 \(\deg(\operatorname{pr}_2)=3\)。
证毕。
\end{proof}

\begin{proposition}[交换商曲线的属 \(2\) 超椭圆模型]\label{prop:conclusion-leyang-symmetric-quotient-genus2}
记交换自同构
\[
\sigma:D\to D,\qquad (P,Q)\mapsto (Q,P),
\]
并令 \(C:=D/\langle \sigma\rangle\)。
取 \(\sigma\)-不变函数
\[
s=\lambda_1+\lambda_2,\qquad p=\lambda_1\lambda_2.
\]
则 \(C\) 与如下超椭圆曲线双有理等价：
\[
z^2=4s^6-12s^5+12s^4-4s^3-4s^2+4s+1,
\]
并且 \(g(C)=2\)。
\end{proposition}
\begin{proof}
由 \(R(\lambda_1,\lambda_2)\) 对称性，\(\sigma\) 正则作用成立。
对方程组
\[
R(\lambda_1,\lambda_2)=0,\qquad
s-(\lambda_1+\lambda_2)=0,\qquad
p-\lambda_1\lambda_2=0
\]
消去 \((\lambda_1,\lambda_2)\)，得
\[
p^2+\bigl(-4s^4+6s^3-2s^2+2\bigr)p+\bigl(s^6-2s^5+s^4+2s^3-3s^2+1\bigr)=0.
\]
令
\[
w:=2p+\bigl(-4s^4+6s^3-2s^2+2\bigr),
\]
则
\[
w^2=4s^2\bigl(4s^6-12s^5+12s^4-4s^3-4s^2+4s+1\bigr).
\]
在开集 \(s\neq 0\) 上令 \(z=w/(2s)\)，即得所示超椭圆方程；故 \(C\) 与其双有理等价。

属数由商覆盖分支计算。
\(\sigma\) 商对应于 \(S_4\) 在无序二元子集 \(\binom{\{1,2,3,4\}}{2}\)（共 \(6\) 点）上的作用，
故 \(C\to\PP^1_y\) 为次数 \(6\) 覆盖。
有限简单分歧换位在该 \(6\)-点集上循环型 \(2^2\,1^2\)，贡献 \(2\)；
无穷远四循环的循环型为 \(4\cdot 2\)，贡献 \((4-1)+(2-1)=4\)。
总分歧度
\[
5\cdot 2+4=14.
\]
由 Riemann--Hurwitz：
\[
2g(C)-2=-2\cdot 6+14=2,
\]
故 \(g(C)=2\)。
证毕。
\end{proof}

\begin{theorem}[有限域纤维计数二阶矩恒等式与有效误差界]\label{thm:conclusion-leyang-finite-field-second-moment}
设 \(S_{\mathrm{bad}}\) 为包含坏约化素数的有限集合，并取 \(p\notin S_{\mathrm{bad}}\)。
定义
\[
N_p(y):=\#\{\lambda\in\mathbb F_p:\ \Pi(\lambda,y)=0\},\qquad y\in\mathbb F_p.
\]
记 \(D_{\mathrm{aff}}\subset D\) 为去除 \(y=\infty\) 纤维后的仿射开曲线。则
\[
\sum_{y\in\mathbb F_p}N_p(y)^2
=\bigl(\#E(\mathbb F_p)-1\bigr)+\#D_{\mathrm{aff}}(\mathbb F_p).
\]
进而存在绝对常数 \(c_0\) 使
\[
\left|
\sum_{y\in\mathbb F_p}N_p(y)^2-2p
\right|
\le 14\sqrt p+c_0,
\]
从而
\[
\frac1p\sum_{y\in\mathbb F_p}N_p(y)^2=2+O(p^{-1/2}).
\]
\end{theorem}
\begin{proof}
对每个有限 \(y\in\mathbb F_p\)，\(N_p(y)\) 是 \(\pi_y^{-1}(y)\cap(E(\mathbb F_p)\setminus\{O\})\) 的基数。
故
\[
\sum_{y\in\mathbb F_p}N_p(y)^2
=\#\{(P,Q)\in(E(\mathbb F_p)\setminus\{O\})^2:\ y(P)=y(Q)\}.
\]
右侧即仿射自纤维积的 \(\mathbb F_p\)-点数，并分解为对角与非对角两部分：
对角贡献 \(\#E(\mathbb F_p)-1\)，非对角贡献 \(\#D_{\mathrm{aff}}(\mathbb F_p)\)。
第一式成立。

再用 Weil--Hasse 界：
\[
\#E(\mathbb F_p)=p+1-a_p(E),\quad |a_p(E)|\le 2\sqrt p,
\]
\[
\#D(\mathbb F_p)=p+1-a_p(D),\quad |a_p(D)|\le 2g(D)\sqrt p=12\sqrt p.
\]
且
\[
\#D_{\mathrm{aff}}(\mathbb F_p)=\#D(\mathbb F_p)-t_{\infty}(p),
\]
其中 \(t_{\infty}(p)\) 为 \(y=\infty\) 纤维上的有理点数，故 \(|1-t_{\infty}(p)|\le c_0\)（取常数上界即可）。
合并得
\[
\left|
\sum_{y\in\mathbb F_p}N_p(y)^2-2p
\right|
\le |a_p(E)|+|a_p(D)|+c_0
\le 14\sqrt p+c_0.
\]
证毕。
\end{proof}

\begin{proposition}[Jacobian 的显式同源分解与四维 Prym 因子]\label{prop:conclusion-leyang-jd-prym-decomposition}
设 \(f_i=\operatorname{pr}_i:D\to E\)（\(i=1,2\)）为定理 \ref{thm:conclusion-leyang-self-fiber-product-genus6} 的次数 \(3\) 投影。
定义
\[
\Phi:J(D)\longrightarrow E\times E,\qquad \Phi=(f_{1*},f_{2*}).
\]
则 \(\Phi\) 为满射，且
\[
P:=\ker(\Phi)^0
\]
是维数 \(4\) 的阿贝尔簇。并存在以 \(3\) 的幂为次数的同源
\[
J(D)\sim_{\QQ}E^2\times P.
\]
\end{proposition}
\begin{proof}
对每个 \(i\)，有限态射 \(f_i\) 满足
\[
f_{i*}\circ f_i^*=[\deg f_i]=[3]\quad\text{于 }E.
\]
故 \(f_{i*}\) 为满射。
由两投影对应的拉回 \(f_1^*,f_2^*:H^0(E,\Omega^1)\to H^0(D,\Omega^1)\) 给出的二维像在 \(H^0(D,\Omega^1)\) 中线性独立，可得 \(\Phi\) 在切空间层面秩为 \(2\)，从而 \(\Phi\) 满射到二维阿贝尔簇 \(E\times E\)。

由 \(g(D)=6\) 得 \(\dim J(D)=6\)，故
\[
\dim P=\dim J(D)-\dim(E\times E)=6-2=4.
\]
Poincar\'e 完全可约性给出
\[
J(D)\sim E^2\times P.
\]
又 \(f_{i*}\circ f_i^*=[3]\) 使对应同源度数仅含素因子 \(3\)，故可取为 \(3\) 的幂次。
证毕。
\end{proof}

\begin{corollary}[属 \(2\) 商模型的分歧判别式]\label{cor:conclusion-leyang-genus2-discriminant}
设
\[
C:\ z^2=f_6(s),\qquad
f_6(s)=4s^6-12s^5+12s^4-4s^3-4s^2+4s+1
\]
为命题 \ref{prop:conclusion-leyang-symmetric-quotient-genus2} 的显式模型。
则
\[
\operatorname{Disc}(f_6)=-2^{12}\cdot 31\cdot 37^2.
\]
因此该模型在素数 \(2,31,37\) 之外不引入新的分歧坏化，并且 \(37\) 以平方指数进入分歧判别式。
\end{corollary}
\begin{proof}
按定义
\[
\operatorname{Disc}(f_6)=(-1)^{15}\operatorname{Res}(f_6,f_6').
\]
直接计算结式并作整数分解即得
\[
\operatorname{Disc}(f_6)=-2^{12}\cdot 31\cdot 37^2.
\]
证毕。
\end{proof}

\begin{remark}[与本质素数谱条目的兼容性]\label{rem:conclusion-leyang-self-fiber-compat-essential-spectrum}
推论 \ref{cor:conclusion-leyang-genus2-discriminant} 中出现的 \(31,37\) 与推论 \ref{cor:conclusion-lee-yang-essential-odd-primes} 的奇素本质谱判据在分歧层面一致；其中 \(37\) 同时作为椭圆坏素数与商曲线判别式平方因子出现，体现了覆盖分支几何与算术坏化的同位锁定。
\end{remark}
